% ------------------------------------------------------------
% Page 107
% ------------------------------------------------------------

Si j’ajoute à cela que les jours de beau soleil il suffisait de sauter la rivière, de rester
immobile à la lisière d’un bosquet de sapins et d’attendre le léger grignotement qui
révélait l’écureuil promis à un délicieux petit civet… Tout cela pour dire que la
viande que n’apportait pas « le Caïffa » ne nous manquait nullement et d’autant plus
que sur les conseils et l’aide de Ricau A. je devins au printemps un pêcheur, disons
passable, mais largement suffisant.

Tous les jours vers dix heures, nous surveillions l’arrivée du facteur, poussant sa
bicyclette, il distribuait le courrier, visitait la boite aux lettres avant d’entrer chez
nous où le café l’attendait. Il déposait à la maison le quotidien de Rousset « La
Dépêche de Toulouse » qu’ils ne consultaient que le soir. Nous avions toute la
journée pour lire et éplucher au jour le jour l’actualité de cette année 1939 fertile en
événements qui furent souvent déterminants pour les années qui suivirent. Le jeune
facteur nous proposait tous les jours ses services pour tout ce dont nous pourrions
avoir besoin. Nous n’abusions pas parce qu’il refusait toujours le prix du morceau de
viande s’estimant lui-même redevable de quelques grives ou plus tard de quelques
truites que nous lui offrions occasionnellement.

\begin{center}
\textbf{Les Habitants des Oubrets}
\end{center}

L’environnement c’était aussi le hameau et ses habitants mais là, je laisse tomber la
plume, car c’est un paragraphe que je devrais consacrer à chacune des familles ;
disons seulement que toutes étaient toujours prêtes à nous être agréables. Ainsi la
brave dame s’arrêtant devant notre porte : « Il fait très froid aujourd’hui, donnez-moi
donc votre petite lessive, je la rincerai avec la mienne ». Il est vrai que ma jeune
épouse avait conquis tout le monde car elle avait su tout de suite visiter mais aussi
recevoir. On venait souvent la voir pour des raisons les plus diverses ( relever un
modèle de tricot, par exemple, ou, tout simplement bavarder).

Bien sûr nous étions invités à toutes les agapes et les festins hivernaux et le
lendemain figurait encore sur notre table la petite « padelade » traditionnelle
(boudins, tranches de rôti, et bout de saucisse).

Et puis encore de temps à autre : « Voici un verre de pâté ou une assiette de cervelas,
j’ai eu peut-être la main un peu lourde pour le sel… on le mangera comme ça
etc..etc.. Nous étions un peu gênés par tant de générosité, nous qui n’avions à
offrir que quelques tasses de café.

Je ne résiste pas à l’envie d’écrire quelques mots sur deux personnes fort différentes
mais faisant preuve à mon égard d’une grande confiance :

\textbf{« POLICE ».} On ne l’appelait qu’ainsi, était un vieux célibataire, un grand-oncle de
Mme. Argeliers --mon ancienne restauratrice du mois d’octobre-- qui l’hébergeait. Il
venait de temps à autre à la maison pour se chauffer peut-être, mais aussi pour me
répéter : » Vous savez, je ne suis pas « un sans le sou », et d’étaler sur la table et de
compter devant moi un pécule de 1450 F. (1939), visiblement heureux devant tant
d’argent qu’il repliait soigneusement avant d’introduire le précieux portefeuille dans
une poche intérieure méticuleusement épinglée.
